% 型8：解法パターン型
% 解法を「型」として番号付きで並べ，手順と適用例をそろえる．
% あとから「型3で処理する」と参照できるため，講座全体の共通語彙になる．
\documentclass[lualatex,paper=b5j,fontsize=10pt,twoside]{jlreq}

\usepackage{surikumi-mondaishu}

\MathMagazineSetup{
  feature={},
  series={解法整理／数学 B},
  series-position=left,
  pattern=pinwheels,
  title={数列の和を求める四つの型},
  author={入試基礎講座},
  author-font=mincho,
  author-position=right,
  title-fill=black,
  deck={和の問題は，どの型に乗せるかを決めた時点で半分終わる．型番号は以降の講義でも使う．}
}

\begin{document}

\MakeMathMagazineTitle
\MSPatternReset

\MSPattern{公式に直接あてはめる}

等差数列と等比数列，および $\sum k$，$\sum k^2$，$\sum k^3$ は公式で処理する．

\begin{msprocedure}
  \item 一般項が $k$ の多項式かどうかを見る．
  \item 多項式なら次数ごとに分けて公式を適用する．
\end{msprocedure}

たとえば
\begin{align*}
  \sum_{k=1}^{n}(2k^2-3k+1)
    &= 2\sum_{k=1}^{n}k^2-3\sum_{k=1}^{n}k+\sum_{k=1}^{n}1\\
    &= \dfrac{n(n+1)(2n+1)}{3}-\dfrac{3n(n+1)}{2}+n
\end{align*}
となる．

\MSPattern{差の形にして途中を消す}

一般項を $b_{k+1}-b_k$ の形に書き直せれば，和は両端だけが残る．

\begin{msprocedure}
  \item 分母が積の形なら部分分数に分解する．
  \item $\sum_{k=1}^{n}(b_{k+1}-b_k)=b_{n+1}-b_1$ を使う．
\end{msprocedure}

たとえば $\dfrac{1}{k(k+1)}=\dfrac{1}{k}-\dfrac{1}{k+1}$ であるから
\[
  \sum_{k=1}^{n}\dfrac{1}{k(k+1)}=1-\dfrac{1}{n+1}=\dfrac{n}{n+1}
\]
となる．

\MSPattern{等比の公比を掛けて引く}

一般項が「$k$ の多項式」×「等比」の形のときに使う．

\begin{msprocedure}
  \item 和を $S$ とおく．
  \item $S$ に公比を掛けた式を作り，$S$ との差をとる．
  \item 残った等比数列の和を計算する．
\end{msprocedure}

$S=\sum_{k=1}^{n}k\cdot 2^{k}$ とおくと $2S-S=S$ の計算から
\[
  S=(n-1)\cdot 2^{n+1}+2
\]
が得られる．

\MSPattern{和を先に予想して帰納法で示す}

規則が見えるが変形の手がかりがないときに使う．

\begin{msprocedure}
  \item $n=1,2,3$ を計算して和の形を予想する．
  \item 数学的帰納法で予想を証明する．
\end{msprocedure}

\MMNote{型3 と 型4 は互いに置き換えられることがある．計算量が少ないほうを選べばよい．}

\MMSubheading{型の判別}

次の和は，上のどの型で処理するのが最も速いか．

\begin{msdrill}
  \item $\displaystyle\sum_{k=1}^{n}(k^2+k)$
  \item $\displaystyle\sum_{k=1}^{n}\dfrac{1}{(k+1)(k+2)}$
  \item $\displaystyle\sum_{k=1}^{n}k\cdot 3^{k}$
  \item $\displaystyle\sum_{k=1}^{n}\dfrac{1}{\sqrt{k+1}+\sqrt{k}}$
\end{msdrill}

\begin{msanswerstrip}
  \MSAns{1}{型1}
  \MSAns{2}{型2}
  \MSAns{3}{型3}
  \MSAns{4}{型2（分母を有理化する）}
\end{msanswerstrip}

\end{document}
